% GENERATED FILE. Edit canonical lessons, then run npm run generate:books.
% canonical-source-hash: fnv1a64:565b6779cc920c98
% canonical-lessons: TE-C463-rayabara-karyalayam, TE-C463-rayabari, TE-C463-ketayincu, TE-C463-podigincu, TE-C463-punaruddharincu

\chapter{The Embassy}
\label{ch:te-the-embassy}
\hlchaptermodality{463}

\begin{chapteropening}
\textbf{By the end of this chapter:} \emph{I can say an embassy, an ambassador, to allot, to extend and to renew.}

Complete the last lesson of chapter 463: I can say an embassy, an ambassador, to allot, to extend and to renew.
\end{chapteropening}

\section[\texorpdfstring{\te{రాయబార} \te{కార్యాలయం} (rāyabāra kāryālayaṁ)}{రాయబార కార్యాలయం (rāyabāra kāryālayaṁ)}]{\te{రాయబార} \te{కార్యాలయం} (rāyabāra kāryālayaṁ) — an embassy}

\label{lesson:TE-C463-rayabara-karyalayam}



\begin{quote}
Before the new one: say the Telugu for to issue, then the Telugu for to sanction.
\end{quote}

\subsection*{You'll want to know: \te{రాయబార} \te{కార్యాలయం}}
\textbf{\te{రాయబార} \te{కార్యాలయం}} — \emph{rāyabāra kāryālayaṁ} — \textquotedblleft{}an embassy\textquotedblright{}.

\begin{quote}
\te{వీసా} \te{కోసం} \te{రాయబార} \te{కార్యాలయానికి} \te{వెళ్ళాను}.
\end{quote}

\begin{quote}
\emph{vīsā kōsaṁ rāyabāra kāryālayāniki veḷḷānu.}
\end{quote}

\begin{quote}
I went to the embassy for a visa.
\end{quote}

\te{రాయబారం}, an embassy or a mission, + \te{కార్యాలయం}, an office.

\begin{cousinweb}
\textbf{\te{రాయబార} \te{కార్యాలయం}} (\emph{rāyabāra kāryālayaṁ}) is a Telugu phrase: \te{రాయబార}, \textquotedblleft{}of an envoy\textquotedblright{}, whose origin is not settled, and \te{కార్యాలయం}, \textquotedblleft{}an office\textquotedblright{}, built from Sanskrit \emph{kārya}, \textquotedblleft{}work\textquotedblright{}, and \emph{ālaya}, \textquotedblleft{}a house\textquotedblright{} (Monier-Williams).
\end{cousinweb}

\begin{grammarlens}[title={What you've built}]
One more word for getting around.
\end{grammarlens}

\subsection*{Your turn}
\begin{itemize}
\raggedright
  \item \emph{Say it:} \emph{rāyabāra kāryālayaṁ}
  \item \emph{Say it:} \emph{rāyabāra kāryālayaṁ}, once more
  \item \emph{Say it:} \emph{vīsā kōsaṁ rāyabāra kāryālayāniki veḷḷānu.} — I went to the embassy for a visa.
  \item \emph{Recall:} say the Telugu for to issue, then the Telugu for to sanction, then say \emph{rāyabāra kāryālayaṁ} again
\end{itemize}

\begin{tcolorbox}[breakable,colback=teal!4,colframe=teal!35!black,title={Before you move on}]
What does \te{రాయబార} \te{కార్యాలయం} mean? (\textquotedblleft{}An embassy\textquotedblright{}.) Say it once more.
\end{tcolorbox}
\section[\texorpdfstring{\te{రాయబారి} (rāyabāri)}{రాయబారి (rāyabāri)}]{\te{రాయబారి} (rāyabāri) — an ambassador; an envoy}

\label{lesson:TE-C463-rayabari}



\begin{quote}
Before the new one: say the Telugu for to sanction, then the Telugu for an embassy.
\end{quote}

\subsection*{You'll want to know: \te{రాయబారి}}
\textbf{\te{రాయబారి}} — \emph{rāyabāri} — \textquotedblleft{}an ambassador; an envoy\textquotedblright{}.

\begin{quote}
\te{మన} \te{దేశ} \te{రాయబారి} \te{ఆ} \te{సభలో} \te{మాట్లాడారు}.
\end{quote}

\begin{quote}
\emph{mana dēśa rāyabāri ā sabhalō māṭlāḍāru.}
\end{quote}

\begin{quote}
Our country's ambassador spoke at that meeting.
\end{quote}

In the old stories a \te{రాయబారి} carried a king's message to another court.

\begin{cousinweb}
\textbf{\te{రాయబారి}} (\emph{rāyabāri}) is a Telugu word for an envoy whose origin is not settled.
\end{cousinweb}

\begin{grammarlens}[title={What you've built}]
One more word for getting around.
\end{grammarlens}

\subsection*{Your turn}
\begin{itemize}
\raggedright
  \item \emph{Say it:} \emph{rāyabāri}
  \item \emph{Say it:} \emph{rāyabāri}, once more
  \item \emph{Say it:} \emph{mana dēśa rāyabāri ā sabhalō māṭlāḍāru.} — Our country's ambassador spoke at that meeting.
  \item \emph{Recall:} say the Telugu for to sanction, then the Telugu for an embassy, then say \emph{rāyabāri} again
\end{itemize}

\begin{tcolorbox}[breakable,colback=teal!4,colframe=teal!35!black,title={Before you move on}]
What does \te{రాయబారి} mean? (\textquotedblleft{}An ambassador; an envoy\textquotedblright{}.) Say it once more.
\end{tcolorbox}
\section[\texorpdfstring{\te{కేటాయించు} (kēṭāyiñcu)}{కేటాయించు (kēṭāyiñcu)}]{\te{కేటాయించు} (kēṭāyiñcu) — to allot; to assign}

\label{lesson:TE-C463-ketayincu}



\begin{quote}
Before the new one: say the Telugu for an embassy, then the Telugu for an ambassador.
\end{quote}

\subsection*{You'll want to know: \te{కేటాయించు}}
\textbf{\te{కేటాయించు}} — \emph{kēṭāyiñcu} — \textquotedblleft{}to allot; to assign\textquotedblright{}.

\begin{quote}
\te{హోటలు} \te{వాళ్ళు} \te{మాకు} \te{పై} \te{అంతస్తులో} \te{గది} \te{కేటాయించారు}.
\end{quote}

\begin{quote}
\emph{hōṭalu vāḷḷu māku pai antastulō gadi kēṭāyiñcāru.}
\end{quote}

\begin{quote}
The hotel people gave us a room on the top floor.
\end{quote}

Past (I …): \te{కేటాయించాను}, \emph{kēṭāyiñcānu}.

Said of rooms, seats, money or duties given out by an office.

\begin{cousinweb}
\textbf{\te{కేటాయించు}} (\emph{kēṭāyiñcu}) is a Telugu verb in -\te{ఇంచు} whose origin is not settled.
\end{cousinweb}

\begin{grammarlens}[title={What you've built}]
One more word for getting around.
\end{grammarlens}

\subsection*{Your turn}
\begin{itemize}
\raggedright
  \item \emph{Say it:} \emph{kēṭāyiñcu}
  \item \emph{Say it:} \emph{kēṭāyiñcu}, once more
  \item \emph{Say it:} \emph{hōṭalu vāḷḷu māku pai antastulō gadi kēṭāyiñcāru.} — The hotel people gave us a room on the top floor.
  \item \emph{Recall:} say the Telugu for an embassy, then the Telugu for an ambassador, then say \emph{kēṭāyiñcu} again
\end{itemize}

\begin{tcolorbox}[breakable,colback=teal!4,colframe=teal!35!black,title={Before you move on}]
What does \te{కేటాయించు} mean? (\textquotedblleft{}To allot; to assign\textquotedblright{}.) Say it once more.
\end{tcolorbox}
\section[\texorpdfstring{\te{పొడిగించు} (poḍigiñcu)}{పొడిగించు (poḍigiñcu)}]{\te{పొడిగించు} (poḍigiñcu) — to extend; to lengthen}

\label{lesson:TE-C463-podigincu}



\begin{quote}
Before the new one: say the Telugu for an ambassador, then the Telugu for to allot.
\end{quote}

\subsection*{You'll want to know: \te{పొడిగించు}}
\textbf{\te{పొడిగించు}} — \emph{poḍigiñcu} — \textquotedblleft{}to extend; to lengthen\textquotedblright{}.

\begin{quote}
\te{రాయబార} \te{కార్యాలయం} \te{నా} \te{వీసాను} \te{ఆరు} \te{నెలలు} \te{పొడిగించింది}.
\end{quote}

\begin{quote}
\emph{rāyabāra kāryālayaṁ nā vīsānu āru nelalu poḍigiñcindi.}
\end{quote}

\begin{quote}
The embassy extended my visa by six months.
\end{quote}

Past (I …): \te{పొడిగించాను}, \emph{poḍigiñcānu}.

Said of time and of length: a visa, a deadline, a road.

\begin{cousinweb}
\textbf{\te{పొడిగించు}} (\emph{poḍigiñcu}) is native Dravidian: DEDR 4484 lists Telugu \emph{poḍigincu}, \textquotedblleft{}to lengthen, to raise\textquotedblright{}, beside \emph{poḍugu}, \textquotedblleft{}length, height\textquotedblright{}.
\end{cousinweb}

\begin{grammarlens}[title={What you've built}]
One more word for getting around.
\end{grammarlens}

\subsection*{Your turn}
\begin{itemize}
\raggedright
  \item \emph{Say it:} \emph{poḍigiñcu}
  \item \emph{Say it:} \emph{poḍigiñcu}, once more
  \item \emph{Say it:} \emph{rāyabāra kāryālayaṁ nā vīsānu āru nelalu poḍigiñcindi.} — The embassy extended my visa by six months.
  \item \emph{Recall:} say the Telugu for an ambassador, then the Telugu for to allot, then say \emph{poḍigiñcu} again
\end{itemize}

\begin{tcolorbox}[breakable,colback=teal!4,colframe=teal!35!black,title={Before you move on}]
What does \te{పొడిగించు} mean? (\textquotedblleft{}To extend; to lengthen\textquotedblright{}.) Say it once more.
\end{tcolorbox}
\section[\texorpdfstring{\te{పునరుద్ధరించు} (punaruddhariñcu)}{పునరుద్ధరించు (punaruddhariñcu)}]{\te{పునరుద్ధరించు} (punaruddhariñcu) — to renew; to restore}

\label{lesson:TE-C463-punaruddharincu}



\begin{quote}
Before the new one: say the Telugu for to allot, then the Telugu for to extend.
\end{quote}

\subsection*{You'll want to know: \te{పునరుద్ధరించు}}
\textbf{\te{పునరుద్ధరించు}} — \emph{punaruddhariñcu} — \textquotedblleft{}to renew; to restore\textquotedblright{}.

\begin{quote}
\te{లైసెన్సును} \te{ప్రతి} \te{ఐదు} \te{సంవత్సరాలకు} \te{పునరుద్ధరించాలి}.
\end{quote}

\begin{quote}
\emph{laisensunu prati aidu saṁvatsarālaku punaruddhariñcāli.}
\end{quote}

\begin{quote}
The licence must be renewed every five years.
\end{quote}

Past (I …): \te{పునరుద్ధరించాను}, \emph{punaruddhariñcānu}.

\te{పునర్}, again, + \te{ఉద్ధరించు}, to raise up: to bring back to life. Used of a licence and of an old building.

\begin{cousinweb}
\textbf{\te{పునరుద్ధరించు}} (\emph{punaruddhariñcu}) is built from Sanskrit \emph{punar}, \textquotedblleft{}again\textquotedblright{}, and \emph{ud-dhṛ}, \textquotedblleft{}to raise up, to rescue\textquotedblright{}, whose noun is \emph{punaruddhāra}, \textquotedblleft{}restoration\textquotedblright{} (Monier-Williams), with the Telugu verb ending -\te{ఇంచు}.
\end{cousinweb}

\begin{grammarlens}[title={What you've built}]
One more word for getting around.
\end{grammarlens}

\subsection*{Your turn}
\begin{itemize}
\raggedright
  \item \emph{Say it:} \emph{punaruddhariñcu}
  \item \emph{Say it:} \emph{punaruddhariñcu}, once more
  \item \emph{Say it:} \emph{laisensunu prati aidu saṁvatsarālaku punaruddhariñcāli.} — The licence must be renewed every five years.
  \item \emph{Recall:} say the Telugu for to allot, then the Telugu for to extend, then say \emph{punaruddhariñcu} again
\end{itemize}

\begin{tcolorbox}[breakable,colback=teal!4,colframe=teal!35!black,title={Before you move on}]
What does \te{పునరుద్ధరించు} mean? (\textquotedblleft{}To renew; to restore\textquotedblright{}.) Say it once more.
\end{tcolorbox}
